\documentclass[12pt]{article}
\usepackage[a4paper,margin=1in]{geometry}
\usepackage[utf8]{inputenc}
\usepackage[T1]{fontenc}
\usepackage{amsmath,amssymb}
\usepackage{textcomp}
\usepackage{parskip}
\begin{document}
\begin{center}
{\Large \textbf{2015 STEP III}}
\end{center}

\subsection*{Page 1}
\begin{flushleft}
\bigskip\textbf{Section A: Pure Mathematics}\par
\bigskip\textbf{Question 1.} (i) Let\par
\infty\par
\[In =\]
\[\leq0 (1 + u^2)n\]
du ,\par
where n is a positive integer. Show that\par
\[In In+1 = In\]
- 2n\par
and deduce that\par
\[(2n)! \pi\]
\[In+1 =\]
22n+1(n!)2\par
.\par
\textbf{(ii)} Let\par
\infty\par
-1 2\par
\[J = f (x x ) dx ,\]
\[\leq0 -\]
\[\geq (\]
where f is any function for which the integral exists. Show that\par
\infty \infty \infty\par
-2 -1 2 -2 -1 2 2\par
\[J = x f (x x ) dx = (1 + x )f (x x ) dx = f u du .\]
\[\leq0 - 2 \leq0 - \leq0\]
\[\geq ( \geq ( \geq (\]
\textbf{(iii)} Hence evaluate\par
\[\infty x2n-2\]
dx ,\par
\[\leq0 (x^4 x^2 + 1)n\]
-\par
where n is a positive integer.\par
\end{flushleft}
\subsection*{Page 2}
\begin{flushleft}
\[2 If s^1, s^2, s^3, . . . and t^1, t^2, t^3, . . . are sequences of positive numbers, we write\]
\[(sn) \leq (tn)\]
to mean\par
\["there exists a positive integer m such that sn \leq tn whenever n \geq m".\]
Determine whether each of the following statements is true or false. In the case of a true\par
statement, you should give a proof which includes an explicit determination of an appropriate\par
m; in the case of a false statement, you should give a counterexample.\par
\textbf{(i)} (1000n) \leq (n2) .\par
\textbf{(ii)} If it is not the case that (sn) \leq (tn), then it is the case that (tn) \leq (sn) .\par
\textbf{(iii)} If (sn) \leq (tn) and (tn) \leq (un), then (sn) \leq (un) .\par
\textbf{(iv)} (n2) \leq (2 n ) .\par
\bigskip\textbf{Question 3.} In this question, r and θ are polar coordinates with r \geq 0 and \pi < θ \leq \pi, and a and b are\par
-\par
positive constants.\par
Let L be a fixed line and let A be a fixed point not lying on L. Then the locus of points that\par
are a fixed distance (call it d) from L measured along lines through A is called a conchoid of\par
Nicomedes.\par
\textbf{(i)} Show that if\par
\[r a sec θ = b , ( )\]
| - | ∗\par
\[where a > b, then sec θ > 0. Show that all points with coordinates satisfying ( ) lie on\]
∗\par
a certain conchoid of Nicomedes (you should identify L, d and A). Sketch the locus of\par
these points.\par
\textbf{(ii)} In the case a < b, sketch the curve (including the loop for which sec θ < 0) given by\par
\[r a sec θ = b .\]
| - |\par
\[Find the area of the loop in the case a = 1 and b = 2.\]
\[[Note: sec θ dθ = ln sec θ + tan θ + C .]\]
| |\par
)\par
\end{flushleft}
\subsection*{Page 3}
\begin{flushleft}
4 (i) If a, b and c are all real, show that the equation\par
3 2\par
\[z + az + bz + c = 0 ( )\]
∗\par
has at least one real root.\par
\textbf{(ii)} Let\par
2 2 2 3 3 3\par
\[S^1 = z^1 + z^2 + z^3, S^2 = z\]
+ z\par
+ z\par
\[, S^3 = z\]
+ z\par
+ z\par
,\par
\[where z^1, z^2 and z^3 are the roots of the equation ( ). Express a and b in terms of S^1\]
∗\par
\[and S^2, and show that\]
\[6c = S\]
\[+ 3S1S2 2S3 .\]
- -\par
\textbf{(iii)} The six real numbers rk and θk (k = 1, 2, 3), where rk > 0 and \pi < θk < \pi, satisfy\par
-\par
3 3 3\par
\[rk sin(θk) = 0 , r\]
k\par
\[2 sin(2θk) = 0 , r\]
k\par
\[3 sin(3θk) = 0 .\]
\[k=1 k=1 k=1\]
\[\int \int \int\]
\[Show that θk = 0 for at least one value of k.\]
\[Show further that if θ1 = 0 then θ2 = θ3 .\]
-\par
5 (i) In the following argument to show that \sqrt2 is irrational, give proofs appropriate for\par
steps 3, 5 and 6.\par
\[1. Assume that \sqrt\{2\} is rational.\]
2. Define the set S to be the set of positive integers with the following property:\par
\[n is in S if and only if n\sqrt\{2\} is an integer.\]
3. Show that the set S contains at least one positive integer.\par
4. Define the integer k to be the smallest positive integer in S.\par
\[5. Show that (\sqrt\{2\} 1)k is in S.\]
-\par
\[6. Show that steps 4 and 5 are contradictory and hence that \sqrt\{2\} is irrational.\]
1 2\par
\textbf{(ii)} Prove that 2 3 is rational if and only if 2 3 is rational.\par
1 2\par
Use an argument similar to that of part (i) to prove that 2 3 and 2 3 are irrational.\par
\end{flushleft}
\subsection*{Page 4}
\begin{flushleft}
\bigskip\textbf{Question 6.} (i) Let w and z be complex numbers, and let u = w + z and v = w2 + z2. Prove that w\par
\[and z are real if and only if u and v are real and u^2 \leq 2v.\]
\textbf{(ii)} The complex numbers u, w and z satisfy the equations\par
\[w + z u = 0\]
-\par
2 2 2 2\par
\[w + z u =\]
- - 3\par
3 3\par
\[w + z λu = λ\]
- -\par
where λ is a positive real number. Show that for all values of λ except one (which you\par
should find) there are three possible values of u, all real.\par
Are w and z necessarily real? Give a proof or counterexample.\par
\bigskip\textbf{Question 7.} An operator D is defined, for any function f, by\par
df(x)\par
\[Df(x) = x .\]
dx\par
n\par
The notation D means that D is applied n times; for example\par
d df(x)\par
\[D f(x) = x x .\]
\[dx \sum dx (\]
\[Show that, for any constant a, D2x a = a2x a .\]
\textbf{(i)} Show that if P(x) is a polynomial of degree r (where r \geq 1) then, for any positive\par
n\par
integer n, D P(x) is also a polynomial of degree r.\par
n m\par
\textbf{(ii)} Show that if n and m are positive integers with n < m, then D (1 x) is divisible\par
m-n -\par
by (1 x) .\par
-\par
\textbf{(iii)} Deduce that, if m and n are positive integers with n < m, then\par
m\par
m\par
r n\par
\[( 1) r = 0 .\]
\[- \sum r (\]
\[)r=0\]
\end{flushleft}
\subsection*{Page 5}
\begin{flushleft}
\bigskip\textbf{Question 8.} (i) Show that under the changes of variable x = r cos θ and y = r sin θ, where r is a\par
\[function of θ with r > 0, the differential equation\]
dy\par
\[(y + x) = y x\]
dx -\par
becomes\par
dr\par
\[+ r = 0 .\]
dθ\par
Sketch a solution in the x-y plane.\par
\textbf{(ii)} Show that the solutions of\par
dy\par
2 2 2 2\par
\[y + x x(x + y ) = y x y(x + y )\]
- dx - -\par
\[\geq (\]
can be written in the form\par
\[r =\]
\[1 + Ae2θ\]
and sketch the different forms of solution that arise according to the value of A.\par
\end{flushleft}
\subsection*{Page 6}
\begin{flushleft}
\bigskip\textbf{Section B: Mechanics}\par
\bigskip\textbf{Question 9.} A particle P of mass m moves on a smooth fixed straight horizontal rail and is attached to\par
a fixed peg Q by a light elastic string of natural length a and modulus λ. The peg Q is a\par
\[distance a from the rail. Initially P is at rest with P Q = a.\]
An impulse imparts to P a speed v along the rail. Let x be the displacement at time t of P\par
from its initial position. Obtain the equation\par
\[x˙ 2 = v 2 k 2 x^2 + a^2 a\]
- -\par
[] \{\par
\[where k^2 = λ/(ma), k > 0 and the dot denotes differentiation with respect to t.\]
Find, in terms of k, a and v, the greatest value, x0, attained by x. Find also the acceleration\par
\[of P at x = x^0.\]
Obtain, in the form of an integral, an expression for the period of the motion. Show that, in\par
the case v ka (that is, v is much less than ka), this is approximately\par
(\par
32a 1\par
du .\par
\[\} kv \leq0 \sqrt\{1\} u^4\]
-\par
\bigskip\textbf{Question 10.} A light rod of length 2a has a particle of mass m attached to each end and it moves in a\par
vertical plane. The midpoint of the rod has coordinates (x, y), where the x-axis is horizontal\par
(within the plane of motion) and y is the height above a horizontal table. Initially, the rod\par
is vertical, and at time t later it is inclined at an angle θ to the vertical.\par
Show that the velocity of one particle can be written in the form\par
x˙ + aθ˙ cos θ\par
\[\sumy˙ aθ˙ sin θ (\]
-\par
and that\par
x¨ + aθ¨cos θ aθ˙2 sin θ sin θ 0\par
\[m - = T mg\]
\[\sumy¨ aθ¨sin θ aθ˙2 cos θ( - \sumcos θ( - \sum1(\]
- -\par
where the dots denote differentiation with respect to time t and T is the tension in the rod.\par
Obtain the corresponding equations for the other particle.\par
\[Deduce that x¨ = 0, y¨ = g and θ¨ = 0.\]
-\par
Initially, the midpoint of the rod is a height h above the table, the velocity of the higher\par
u 0\par
particle is , and the velocity of the lower particle is . Given that the two particles\par
v v\par
[ \{ [ \{\par
hit the table for the first time simultaneously, when the rod has rotated by \pi, show that\par
2 2 2\par
\[2hu = \pi a g 2\piuva .\]
-\par
\end{flushleft}
\subsection*{Page 7}
\begin{flushleft}
\bigskip\textbf{Question 11.} (i) A horizontal disc of radius r rotates about a vertical axis through its centre with\par
angular speed ω. One end of a light rod is fixed by a smooth hinge to the edge of\par
the disc so that it can rotate freely in a vertical plane through the centre of the disc.\par
A particle P of mass m is attached to the rod at a distance d from the hinge. The\par
rod makes a constant angle α with the upward vertical, as shown in the diagram, and\par
\[d sin α < r.\]
P\par
d\par
α\par
By considering moments about the hinge for the (light) rod, show that the force exerted\par
on the rod by P is parallel to the rod.\par
Show also that\par
\[r cot α = a + d cos α ,\]
g\par
\[where a = . State clearly the direction of the force exerted by the hinge on the rod,\]
ω2\par
and find an expression for its magnitude in terms of m, g and α.\par
\textbf{(ii)} The disc and rod rotate as in part (i), but two particles (instead of P ) are attached to\par
\[the rod. The masses of the particles are m^1 and m^2 and they are attached to the rod\]
\[at distances d^1 and d^2 from the hinge, respectively. The rod makes a constant angle β\]
\[with the upward vertical and d^1 sin β < d^2 sin β < r. Show that β satisfies an equation\]
of the form\par
\[r cot β = a + b cos β ,\]
\[where b should be expressed in terms of d^1, d^2, m^1 and m^2.\]
\end{flushleft}
\subsection*{Page 8}
\begin{flushleft}
\bigskip\textbf{Section C: Probability and Statistics}\par
\bigskip\textbf{Question 12.} A 6-sided fair die has the numbers 1, 2, 3, 4, 5, 6 on its faces. The die is thrown n times,\par
the outcome (the number on the top face) of each throw being independent of the outcome\par
of any other throw. The random variable Sn is the sum of the outcomes.\par
\textbf{(i)} The random variable Rn is the remainder when Sn is divided by 6. Write down the\par
\[probability generating function, G(x), of R^1 and show that the probability generating\]
\[function of R^2 is also G(x). Use a generating function to find the probability that Sn\]
is divisible by 6.\par
\textbf{(ii)} The random variable Tn is the remainder when Sn is divided by 5. Write down the\par
\[probability generating function, G^1(x), of T^1 and show that G^2(x), the probability\]
\[generating function of T^2, is given by\]
1 2\par
\[G^2(x) =\]
(x + 7y)\par
\[where y = 1 + x + x^2 + x^3 + x^4 .\]
Obtain the probability generating function of Tn and hence show that the probability\par
that Sn is divisible by 5 is\par
1 1\par
\[5 \sum - 6n (\]
if n is not divisible by 5. What is the corresponding probability if n is divisible by 5?\par
13 Each of the two independent random variables X and Y is uniformly distributed on the\par
interval [0, 1].\par
\textbf{(i)} By considering the lines x + y = constant in the x-y plane, find the cumulative distri-\par
bution function of X + Y .\par
Hence show that the probability density function f of (X + Y )\par
-1\par
is given by\par
\[2t -2 t -3 for 1 \leq t \leq 1\]
- 2\par
\[f(t) = ⎧ t -3 for 1 \leq t <\]
⎪ ⎨ \infty\par
0 otherwise.\par
⎪\par
⎩\par
Evaluate E .\par
X + Y\par
[ \{\par
\textbf{(ii)} Find the cumulative distribution function of Y /X and use this result to find the prob-\par
X\par
ability density function of .\par
X + Y\par
X\par
Write down E and verify your result by integration.\par
X + Y\par
[ \{\par
\end{flushleft}
\end{document}
